This class extends ArrayList. A TDE list is the intermediate form of generated
code. The list has variable declarations at the front followed by logic using
those variables. An index, 'variables', is used to maintain the insertion point
for variables at the head of the list. Method addTDE() appends code TDEs to the
tail of the list. Method addVTDE inserts TDEs at the end of the leading variable
section of the list. Field 'index' provides a unique integer to append to the
end of generated signal names. It is incremented after each use.

There are a number of additional static lists of particular TDEs that are used
for optimisation. These are -

\blist
\item connects - CONNECT TDE
\item dels - DEL TDE
\item regs - REG TDE
\item ops - OPERATOR TDE
\item whens - WHEN TDE
\item iloops - ILOOP TDE
\item waits - WAIT TDE
\item gates - AND, OR, XOR TDE
\item sels - SELECT TDE
\item execps - EXECP TDE
\item dffs - DFF TDE
\blend

These lists avoid the need to scan the entire TDE list when optimising
particular TDE types.

Another static list, {\bf branches}, contains details of branch statements for
later optimisation.

The static list {\bf start\_dffs} is built up as ILOOPs are optimised. It is
later scanned to eliminate duplicate DFF elements with common start signal.

{\bf static  HashMap\verb'<'String,TDE\verb'>' uelements} is a map of
elements created using the ELEMENT TDE. Static {\bf current\_uelement} is
the most recent ELEMENT TDE to which connections can be added without
explicitly using the element identifier.



\subsection{method optimise}

    This method is called from the Main method (ThreePL.jj) after
    immediate execution. Optimisation, including that done in this
    method, is described in a later chapter.

\subsection{method ncode}

    This method generates the EDIF netlist. It does this in the following
    steps -
    
    The TDE list is traversed (skipping the initial signal declaration
    section) calling TDE.ncode() on each to generate netlist elements and nets
    via XTDECode methods.
    
    \begin{lstlisting}[gobble=8, language=java]
        Iterator<TDE>   it   = iterator();
        TDE             tde  = null;
        while (it.hasNext()) {
        tde = it.next();
        if (tde.isActive())
            tde.ncode(family);
        }
    \end{lstlisting}
    
    A list of external ports is built from data previously extracted
    when processing IBUF and OBUF TDEs in XTDECode.

    \begin{lstlisting}[gobble=8, language=java]
        family.makePorts();
    \end{lstlisting}
    
    Any flip-flops in output pads whose outputs are used internally as well are
    duplicated by a CLB flip-flop (XNetListOptimise.optimiseOBUF() and
    XNetListOptimise.optimiseOBUFT()). Gates with output not connected, gates
    with one input and gates with duplicated inputs are removed or modified as
    appropriate. Family-specific devices such as XORCY, MUXF etc. are similarly
    optimised.
     
    \begin{lstlisting}[gobble=8, language=java]
        passes = family.optimiseNetlist();
    \end{lstlisting}

    Generic elements are converted to family-specific ones where necessary.
    Up to this point gates are as wide as required by the context. Here they are
    converted to combinatorial logic matching the CLBs in the FPGA family.
        
    \begin{lstlisting}[gobble=8, language=java]
        family.convertGenericToSpecific();
    \end{lstlisting}

    The netlist file header is written out.
    
    The netlist cells are then written out to the netlist file.
    
    \begin{lstlisting}[gobble=8, language=java]
        Element.outputCells(pw, view);
    \end{lstlisting}

    The list of external ports is now written to the netlist file.
    
    \begin{lstlisting}[gobble=8, language=java]
        TDECode.outputExterns(pw);
    \end{lstlisting}

    Element instances are written.
   
    \begin{lstlisting}[gobble=8, language=java]
        Element.outputInstances(pw, view);
    \end{lstlisting}

    The nets are now written.
    
    \begin{lstlisting}[gobble=8, language=java]
        Net.outputInstances(pw);
    \end{lstlisting}

    The netlist file is now completed by writing out the part specification
    and the final curly bracket.
    
    Finally constraints are written to the .ncf (ISE) or .xdc (Vivado) file.

    \begin{lstlisting}[gobble=8, language=java]
        family.outputConstraints(author, netlistcomments);
    \end{lstlisting}

